\section{Detection: Evidence and Access}\label{sec:detection}
Detection estimates errors under a specified reference; Figure~\ref{fig:structure} groups methods by evidence and model access.

\subsection{External evidence}
Evidence-based methods compare outputs with supplied or retrieved knowledge. QAFactEval asks questions about summary content and checks answer recovery from the source; question generation and answerability affect its decisions \citep{fabbri-etal-2022-qafacteval}. AlignScore learns alignment across heterogeneous tasks and aggregates source-chunk/output-sentence scores \citep{zha-etal-2023-alignscore}. Both use trained auxiliaries; neither guarantees correct verification.

FActScore decomposes long answers into atomic facts and estimates their supported fraction relative to a designated source \citep{min-etal-2023-factscore}. VeriScore instead extracts verifiable claims, using surrounding context to resolve references \citep{song-etal-2024-veriscore}. Extraction changes the score denominator; missing support differs from contradiction. PrefixNLI extends entailment to incomplete prefixes, enabling feedback during decoding \citep{harary-etal-2026-prefixnli}.

\subsection{Sampling and consistency}
SelfCheckGPT scores a response sentence against additional responses to the same prompt \citep{manakul-etal-2023-selfcheckgpt}. Its NLI and prompting variants aggregate contradiction or non-support judgments across samples. This requires repeated generation rather than target-model probabilities; agreement can still reflect a repeated misconception.

Semantic uncertainty instead clusters sampled answers by mutual entailment and estimates entropy over meanings, discounting paraphrase variation \citep{extsemanticuncertainty2023,extsemanticentropy2024}. The weighted estimator needs sequence probabilities; a discrete approximation uses cluster frequencies. SelfCheckGPT localizes inconsistency in a given response, whereas semantic entropy measures dispersion among answer meanings, targeting unstable confabulations. Both depend on sampling and comparison quality; stable false beliefs need not trigger either signal.

\subsection{Internal signals}
FOCUS combines token probabilities with attention-based uncertainty propagation from an accessible proxy model \citep{zhang-etal-2023-enhancing-uncertainty}. The target generator can remain black-box; the proxy must expose both signals. Hidden-state classifiers learn truth labels from activations \citep{azaria-mitchell-2023-internal}. Lookback Lens uses context-versus-generation attention features with a supervised linear classifier \citep{chuang-etal-2024-lookback}. Probes require representative training labels.

A separable activation pattern need not represent truth: associated falsehoods can resemble factual recall in controlled experiments \citep{cheang-etal-2026-llms}. Probes therefore need evaluation across natural error types and generators.

Source alignment cannot establish source truth; sampling and probes supply no independent factual evidence. ``Black-box'' and ``reference-free'' differ because access to generators, proxies, and verifiers must be specified separately.
